\notepage{65}{Pauli matrices}
In the spin-up/spin-down basis,
\[
 \mathbf S=\frac{\hbar}{2}\boldsymbol\sigma,
 \qquad \{A,B\}\equiv AB+BA.
\]
The following identities are useful:
\begin{enumerate}[label=(\arabic*)]
 \item For $i\ne j$, $\anticomm{\sigma_i}{\sigma_j}=0$, or
       $\sigma_i\sigma_j=-\sigma_j\sigma_i$.
 \item Since $S_xS_y-S_yS_x=\ii\hbar S_z$,
 \[
 \sigma_x\sigma_y-\sigma_y\sigma_x=2\ii\sigma_z.
 \]
 Together with (1), this gives $\sigma_x\sigma_y=\ii\sigma_z$ and the
 cyclic permutations.
 \item $\Tr\sigma_i=0$.
 \item $\sigma_i^2=I$; more generally,
       $(\hat{\mathbf n}\cdot\boldsymbol\sigma)^2=I$ for a unit vector
       $\hat{\mathbf n}$.
\end{enumerate}
\smallheading{Proof of the fourth identity}
Write an arbitrary spin-$\tfrac12$ state as
\[
 \ket{\psi}=\alpha\ket{+}+\beta\ket{-},\qquad
 \ket{\pm}\equiv\ket{\hat{\mathbf n},\pm}.
\]
Then
\begin{align*}
 &\left(\hat{\mathbf n}\cdot\mathbf S+\frac{\hbar}{2}I\right)
  \left(\hat{\mathbf n}\cdot\mathbf S-\frac{\hbar}{2}I\right)\ket{\psi}\\
 &\qquad=\left(\hat{\mathbf n}\cdot\mathbf S+\frac{\hbar}{2}I\right)
      (-\beta\hbar\ket{-})=0.
\end{align*}
Because this holds for every state,
\[
 (\hat{\mathbf n}\cdot\mathbf S)^2=\frac{\hbar^2}{4}I,
 \qquad(\hat{\mathbf n}\cdot\boldsymbol\sigma)^2=I.
\]
\editorial{The intermediate action of the second factor is
$-\beta\hbar\ket{-}$; the handwritten intermediate line omits its minus
sign. The final identity is unchanged.}
Combining these results gives the compact Pauli algebra:
\[
 \boxed{\anticomm{\sigma_i}{\sigma_j}=2\delta_{ij}I},\qquad
 \comm{\sigma_x}{\sigma_y}=2\ii\sigma_z
 \quad\text{and cyclic permutations}.
\]
